\chapter{Tervezés}
\label{ch:tervezes}
\section{Az adathalmazok kiválasztása és előkészítése}
Az előzmények között ismertetett adathalmazok közül a Nutrition5k, az OmniFood8k, a DFood-Tj, az ACETADA, a SNAPMe
és a NutriBench adathalmazokat tudtam beszerezni. A források felkutatásához a DuckDuckGo webes keresőt, 
valamint segédeszközként a Claude Opus 5.5 nagy nyelvi modellt használtam. Azt, hogy a letöltött 
adathalmaz valóban az előzményekben ismertetett publikációhoz tartozik, minden esetben ellenőriztem: 
vagy a letöltési oldal nevezi meg a publikációt (címét, szerzőit, arXiv- vagy DOI-azonosítóját), vagy 
a publikáció hivatkozik a letöltési helyre. Az ASA24-ből csak hiányos változatot találtam, ezért azt nem vizsgáltam 
tovább. A felhasznált adathalmazok letöltési helyét, licencét és a használt változatot 
\az{\ref{tab:adathalmaz_forrasok}}.~táblázat foglalja össze.

\begin{table}[htbp]
\centering
\footnotesize
\caption{A felhasznált adathalmazok letöltési helye, licence és a használt változat.}
\label{tab:adathalmaz_forrasok}
\setlength{\tabcolsep}{5pt}
\begin{tabular}{l>{\raggedright\arraybackslash}p{7.2cm}ll}
\hline
Adathalmaz & Letöltési hely & Licenc & Változat \\
\hline
Nutrition5k & \href{https://github.com/google-research-datasets/Nutrition5k}{github.com/google-research-datasets/Nutrition5k} & CC BY 4.0 & -- \\
OmniFood8k  & \href{https://yudongjian.github.io/OmniFood8K-food/}{yudongjian.github.io/OmniFood8K-food} & nincs megadva & -- \\
DFood-Tj    & \href{https://doi.org/10.57760/sciencedb.28012}{doi:10.57760/sciencedb.28012} & CC0 & V3 \\
ACETADA     & \href{https://skynet.ecn.purdue.edu/~coburn6/ACETADA/}{skynet.ecn.purdue.edu/\textasciitilde coburn6/ACETADA} & CC BY-NC 4.0 & -- \\
SNAPMe      & \href{https://doi.org/10.15482/USDA.ADC/1528346}{doi:10.15482/USDA.ADC/1528346} & CC BY-SA 4.0 & 1. \\
NutriBench  & \href{https://huggingface.co/datasets/dongx1997/NutriBench}{huggingface.co/datasets/dongx1997/NutriBench} & CC BY-NC-SA 4.0 & v1, natural \\
\hline
\end{tabular}
\end{table}

Ezután kialakítottam egy mappaszerkezetet, amiben minden adathalmaz külön mappát kapott, és ezekbe elkészítettem egy \texttt{loader.py} Python programot mindegyik számára,
amelyekben az adatok egységes betöltését és elérhetővé tételét implementáltam. Az adathalmazok fölött a mappaszerkezetben előkészítettem egy \texttt{benchmark.py} nevű Python
programot, amely azért felelős, hogy a tápértékbecslési feladat méréseit végrehajtsa. Az almappából ez a komponens be tudja tölteni az adathalmazok adatait kezelő
függvényeket. Ezeket az adatokat a függvények visszatérési értékeiként kapja meg.

\paragraph{}A SNAPMe adathalmazból a nem csomagolt ételek evés előtti fotóit és az egyadagos csomagolt 
termékek csomagolásának elejéről készült fotóit használom. A csomagolt termékek fotóit nem zártam ki, 
így a döntési mérés 199 tesztelt SNAPMe-ételéből 26 csomagolt termék. Ezeknél a becslés inkább 
a termék felismerésén múlik, mint az adag becslésén.

Az adatok feldolgozásához a Pandas programkönyvtárat használtam, a Pythonnak pedig a 3.14-es verzióját. A modelleket a szolgáltatók hivatalos Python-könyvtárain keresztül hívtam, a \texttt{google-genai} 2.33.0-s, az \texttt{anthropic} 2.1.0-s és az \texttt{openai} 3.37.0-s verziójával, a teljes csomaglistát a \texttt{requirements.txt} tartalmazza.

\section{A mérési környezet kialakítása}
A \texttt{benchmark.py} a tesztméréseket a következő módon valósítja meg:
Megkeresi és betölti a \texttt{loader.py} programokat az alkönyvtárakból.
Betölti a választott adathalmazokat egy Pandas dataframe objektumba. A fájl konfigurációs változói segítségével
lehetőség van az adathalmazok részhalmazát használni. A részhalmazok  kiválasztása egy véletlenszerű sorbarendezés után, vagy állítható tartománnyal történik. A véletlenszerű kiválasztáshoz megadható a
reprodukálhatóság érdekében egy seed érték, valamint a méret meghatározására egy "n" paraméter szolgál, amely megadja, hogy hányadik indexig kell a kiválasztást
végezni. A véletlenszerű sorbarendezést a numpy könyvtár végzi. 
A DFood-Tj-ben minden kép külön rekord, és egy étkezésről két kép készült, ezért ott a 
kiválasztás egysége nem a kép, hanem az étkezés, hogy egy étkezés két képe együtt maradjon, 
és az \texttt{n} paraméter az étkezések számát adja meg.
A program az étkezéseket egy sorszámmal azonosítja, és véletlenszerű kiválasztásnál a seedből és ebből a sorszámból képzett 
szöveg SHA-256 hash-értéke szerint rendezi sorba őket. Így a sorrend véletlenszerű, de minden 
futásban azonos, és egy étkezés helye akkor sem változik, ha a szűrés más étkezéseket kihagy. 
Az állítható tartomány módban a többi adathalmazon a rekordok azonosító (sample\_id) szerint 
rendezett sorrendjében, a DFood-Tj-n pedig az étkezések sorszáma szerint lehet kijelölni a 
használandó résztartományt.
A program ezután fájlba írja a mérés konfigurációs paramétereit.
Betölti az adathalmazhoz tartozó képeket és a konfigurációs változók valamint a megadott argumentumok szerint megformázza a mérni kívánt nagy nyelvi modellek
felé küldendő szöveges prompt utasítást. 
Változók segítségével konfigurálható, hogy mely nagy nyelvi modellek milyen paraméterekkel dolgozzák fel a kéréseket, milyen API végponton.
Ezután elküldi a kérést a program az API végpont felé. Állítható, hogy sikertelen válasz esetén hány újrapróbálkozás megengedett.
A prompt utasítások szerint json formátumban kell a nagy nyelvi modelleknek válaszolniuk a kérésekre, ezt a program a kéréssel elküldött JSON-sémával is előírja, és a következő lépésben ezen formátum szerint történik a szöveges
válasz feldolgozása. Ha ebben hiba történik, abban az esetben már nincs újrapróbálkozás, hanem nem értelmezhetőként naplózásra kerül a kapott válasz.
Miután az adathalmaz minden kiválasztott rekordjához megtörténtek a kérések az API végpont felé, a kapott eredményekről egy összefoglaló csv fájlt ír ki a program.

Minden méréshez megadható kimeneti mappa, egy mérés során tetszőleges számú adathalmazt és modellt lehet mérni, minden megadott adathalmazon minden megadott modell
válaszai mentésre és kiértékelésre kerülnek.

\section{A mért mennyiségek kiválasztása}
\label{sec:mert_mennyisegek}
\paragraph{}A feladatkiírásban megfogalmazottak szerint a megoldások összehasonlításakor a becslés pontossága mellett a válaszidőt, a költséget, a stabilitást és a konzisztenciát is figyelembe kell venni. A következő mennyiségek mérése mellett döntöttem:
\begin{enumerate}
    \item Pontosság: a becsült tömeg, energia, szénhidrát, fehérje és zsír eltérése az
    adathalmaz szerinti referenciaértéktől.
    \item Válaszidő: az API-hívás elküldése és a válasz teljes megérkezése között eltelt idő.
    Újrapróbálkozás esetén csak a sikeres kísérlet ideje számít.
    \item Költség: a hívás bemeneti és kimeneti tokenszámából a szolgáltató által a mérés
    idején közzétett listaárak alapján számított költség, dollárban.
    \item Stabilitás és a strukturált kimenet megbízhatósága: az újrapróbálkozások után is sikertelen
    API-hívások száma, a fel nem dolgozható válaszok száma, valamint azon válaszok száma, amelyek nem teljesen JSON formátumban érkeztek, 
    és csak a JSON-objektum szövegből való kiemelése után váltak feldolgozhatóvá.
    \item Konzisztencia: ugyanarra a bemenetre, azonos beállítással ismételt hívások
    becsléseinek eltérése. Ételenként az ismételt becslések szórása, és ezek
    négyzetes átlaga a referenciaértékek átlagához viszonyítva. Az ételenkénti szórások 
    négyzetes átlaga az ismételt mérések ingadozásának bevett mérőszáma, az alanyon belüli 
    szórás \cite{bland1996measurement}; itt az alany az étel. Így az eredmény a pMAE-vel
    azonos skálán értelmezhető.
\end{enumerate}


\section{A vizsgált nagy nyelvi modellek kiválasztása}
\paragraph{}A nagy nyelvi modellek kiválasztásával kapcsolatban a feladatkiírás
a következő követelményeket fogalmazza meg:
A modelleknek elérhetőnek kell lenniük API végponton keresztül,
támogatniuk kell képi és szöveges bemenetet.
Emellett figyelembe kell venni az elérés késleltetését és árazását is a vizsgálat
során.
A becslések gépi feldolgozásához a strukturált kimeneti formátumra van szükség.

\paragraph{}A három vezető szolgáltató, a Google, az OpenAI és az Anthropic
modellkínálata egyaránt három szintre tagolódik: egy olcsó, nagy tokenszám mellett is gazdaságosra tervezett
modellre, egy közepes árú általános modellre és egy csúcsmodellre. A kereskedelmi
modellek felépítése és paraméterszáma nem nyilvános, ezért a modellek képességét nem a
méretük, hanem a szolgáltató által meghatározott termékkategória és ár alapján tudom meghatározni.
Szolgáltatónként a legolcsóbb és a középső szint vizsgálata mellett döntöttem. Egy étkezésnaplózó
alkalmazás felhasználónként naponta több hívást indít, ezért a működtetés költsége
szempontjából releváns kérdés, hogy az olcsóbb szint mennyit veszít a pontosságból a
középsőhöz képest.

\paragraph{}A promptok és a mérést meghatározó további paraméterek hatásának vizsgálatát többségében a Gemini Flash-Lite modellen
végeztem. Ennek oka, hogy a vizsgált modellek közül ez érhető el a legkisebb költséggel, és Sterling és társai \cite{sterling2026vlm} is ezt a modellcsaládot vizsgálták 
ugyanazon az adathalmazon és hasonló prompttal. Annak ellenére, hogy a minták kiválasztása eltérő, az eredmények nagyságrendi viszonyítási alapot adnak a mérési környezet ellenőrzéséhez.

\paragraph{}A modellek összehasonlító méréseiben a reasoning-szintet minden olyan 
modellnél alacsonyra állítottam, ahol ez állítható, hogy a szolgáltatók modelljei 
összevethető beállítással fussanak, a válaszidő és a költség pedig korlátozott 
maradjon. A Gemini Flash-Lite modelleknél az alapértelmezett, minimális szintet 
használtam. Egyes előmérésekben magasabb reasoning-szintet is kipróbáltam.
A vizsgált modelleket és beállításaikat a \ref{tab:modellek}.~táblázat foglalja össze.

\begin{table}[h]
\centering
\caption{A vizsgált modellek. Az árak a szolgáltatók által a mérés idején (2026.
október--november) közzétett listaárak, millió bemeneti/kimeneti tokenenként.}
\label{tab:modellek}
\begin{tabular}{llllr}
\hline
Szolgáltató & Szint & API-azonosító & Beállítás & Ár (\$) \\
\hline
Google    & olcsó & \texttt{gemini-3.5-flash-lite} & temperature 0,1 & 0,30 / 2,50 \\
Google    & olcsó & \texttt{gemini-3.1-flash-lite} & temperature 0,1 & 0,25 / 1,50 \\
Google    & közép & \texttt{gemini-3.8-flash} & thinking: low & 0,75 / 3,75 \\
OpenAI    & olcsó & \texttt{gpt-6-luna} & reasoning: low & 0,10 / 0,50 \\
OpenAI    & közép & \texttt{gpt-6.1-sol} & reasoning: low & 2,00 / 10,00 \\
OpenAI    & csúcs & \texttt{gpt-6-astra} & reasoning: low & 10,00 / 50,00 \\
Anthropic & olcsó & \texttt{claude-haiku-5-5} & effort: low & 0,10 / 0,50 \\
Anthropic & közép & \texttt{claude-sonnet-5-5} & effort: low & 2,00 / 10,00 \\
Anthropic & csúcs & \texttt{claude-opus-5-5} & effort: low & 4,00 / 20,00 \\
\hline
\end{tabular}
\end{table}

A Gemini 3.1-Flash-Lite modellt csak a DFood-Tj adathalmazon végzett fejlesztési mérésekben használtam.

A csúcsszintet kiegészítő felső korlátként vizsgáltam. 
Az Anthropic esetében minden mérésre kiválasztott adathalmazon mértem, 
mert ott a csúcsmodell ára a középső szintnek csak a kétszerese. 
Az OpenAI csúcsmodellje a középső szint ötszörösébe kerül, ezért költségkeret miatt csak az OmniFood8k adathalmazon mértem. 
A Google csúcsmodelljét ugyanezen okból nem vizsgáltam.

\section{A vizsgált prompt utasítások összeállítása}
\label{sec:promptok}
\paragraph{}Egy szöveges prompt megadása a nagy nyelvi modellek használata során azt a célt szolgálja,
hogy annak bemeneteként meghatározza annak számára az elvégzendő feladatot.
A méréseimhez a Sterling és társai \cite{sterling2026vlm} által közölt P3 promptból indultam ki,
amely arra kéri a  nagy nyelvi modellt, hogy a fényképen látható tényleges adagot becsülje, és
ne egy feltételezett szokásos adagméretre kerekítsen.
\paragraph{}A vizsgált változatok a következők:
\begin{enumerate}
    \item \textbf{P3:} Sterling és társai \cite{sterling2026vlm} munkájából átvett utasítás. A modell egy étkezés képén több tételt is felsorolhat, tételnek a modell által külön felsorolt ételt, italt vagy összetevőt nevezem. Az utasítás tételenként kéri a tömeg és az energia becslését, az étkezés becslése pedig a tételek becsléseinek összege.
    \item \textbf{P3\_macros:} ugyanaz a szöveg, kiegészítve a szénhidrát-, fehérje- és
    zsírtartalom kérésével tételenként, valamint a kimeneti szerkezet megfelelő
    bővítésével.
    \item \textbf{P3\_macros\_hint:} a P3\_macros szövege, előtte az étkezés tételeinek felsorolásával, amelyet a
SNAPMe referenciarekordjából képeztem: a résztvevő által az étkezési naplóban kiválasztott ételek FNDDS-leírása,
a tömegük nélkül. Ez a szöveg nem a felhasználó szabad leírása, hanem a referenciaannotációból származik.
    \item \textbf{NutriBench:} Csak szöveges bemenetet használó prompt. Előírja, hogy a
    leírás minden ételét és italát külön tételként kell felsorolni, a megadott
    mennyiséget kell használni (a háztartási mértékeket az adott étel jellemző
    tömegére átváltva), és az energiaérték a makrotápanyagokkal összhangban kell hogy legyen.
    \item \textbf{P3\_prep:} a P3\_macros szövege kiegészítve egy utasítással, amely szerint
    minden tételt konkrétan kell megnevezni, az elkészítési módjával és a látható
    hozzáadott összetevőkkel együtt, a látható szószokat, öntetet, olajat, kenhető zsiradékot és italokat pedig
    külön tételként kell felsorolni. Egyetlen hívás, a kimenet azonos a P3\_macros
    kimenetével.
    \item \textbf{P3\_prep\_textcalc:} kétlépéses eljárás. Az első lépés a P3\_prep a
    fényképen. A második lépés csak szöveget használ: az első lépés tételeit
    a NutriBench-hez használt utasításba helyezi, amely a
    megadott tömegekből újra megbecsüli a tápértékeket. A végső előrejelzésben az
    energia és a makrók a második lépésből, a tömeg az első lépésből származik. A válaszidő, a tokenszám és a
    költség a két hívás összege. A kimenet szerkezete itt is azonos a P3\_macros
    kimenetével.
    \item \textbf{P3\_recipe:} a P3\_macros kiegészítése a képen nem látható energiatartalom 
    becslésére, az Open-KNEAD recipe-prior lépéséhez \cite{coburn2026openknead} hasonló elven. 
    A prompt előírja a becslési feladat lépéseit. A modell először megnevezi az ételt, a konyháját és 
    szokásos elkészítési módját, a hozzá tartozó zsír- és cukortartalommal együtt, akkor is, ha 
    az a képen nem látszik. Ezután a kép alapján dönti el, mi és mennyi van a tányéron, az étel 
    szokásos elkészítése alapján pedig azt, hogy grammonként mit tartalmaz. Az elkészítéshez használt 
    olaj, vaj, szósz vagy cukor annak az ételnek az értékeibe számít bele, amellyel együtt készült, külön 
    tételként csak akkor szerepel, ha külön tálalják. A P3\_macros-tól a szénhidrát fogalmának meghatározásában is 
    eltér (összes szénhidrát, a rosttal együtt). A kimeneti szerkezet megegyezik a P3\_macros-éval.
    \item \textbf{P3\_anchor:} a P3\_macros tételenkénti utasítását háromlépéses számításra cseréli. 
    A modell előbb a tétel tömegét becsli, majd emlékezetből megadja a tétel 100 grammra 
    vonatkozó összetételét, ahogy egy tápanyagtáblázatban szerepelne (táblázatot vagy adatbázist 
    nem kap), és ebből számolja a tápanyagokat. Az energiát az általános Atwater-faktorokkal 
    számolja: 4~kcal/g szénhidrát, 4~kcal/g fehérje és 9~kcal/g zsír \cite{fao2003foodenergy}. 
    A prompt ételcsoportonként a jellemző szénhidrát- és zsírtartalmat is megadja 100 grammra, 
    amelyhez a modell az értékeit viszonyíthatja.
    \item \textbf{P3\_portion:} az adag becsléséhez ismert méretű tárgyakat ad meg viszonyítási 
    alapként (tányér, evőeszköz, pohár, üdítősdoboz). A modell a tételek térfogatát becsli, és 
    jellemző sűrűségekkel grammra váltja; a halmozott vagy részben takart ételt is figyelembe veszi, 
    és ha két méret is lehetséges, azt választja, amelyet a viszonyítási tárgyak alátámasztanak, 
    nem a szokásos adagot. A leírásmezőben meg kell neveznie a használt viszonyítási tárgyat.
    \item \textbf{P3\_check:} a modell a leírásmezőben előbb vázlatot ír minden tétel 
    tömegéről, szénhidrát-, fehérje- és zsírtartalmáról és energiájáról, majd a vázlatot négy 
    szempont szerint ellenőrzi: a becsült energia legfeljebb kb. 15\%-kal térhet el attól, ami a 
    makrotápanyagokból az Atwater-faktorokkal adódik (4~kcal/g szénhidrát és fehérje, 9~kcal/g 
    zsír \cite{fao2003foodenergy}); a három makrotápanyag grammban mért összege nem haladhatja 
    meg a tömeget (főtt ételeknél a prompt szerint ez jellemzően a tömeg 10--40\%-a, olajoknál és 
    zsiradékoknál legfeljebb 100\%); minden étel és ital pontosan egyszer szerepelhet; az összes 
    tömegnek és energiának hihetőnek kell lennie a képen látható ételhez. A javításokat a leírásmezőben 
    rögzíti, a javított értékeket a tétellistába írja.
    \item \textbf{P3\_recipe\_check:} a P3\_recipe és a P3\_check egyesítése: előbb az étel 
    szokásos elkészítése, utána a vázlat és az ellenőrzés.
    \item \textbf{P3\_density:} a P3-ból indul, és tételenként a számok előtt a következőket kéri: 
    az étel nevét, elkészítési módját, a látható hozzáadott zsiradék mértékét (nincs, kevés, 
    közepes, sok), a legközelebbi tápanyagtáblázati tételt és az étel energiasűrűségét 
    (kcal/100~g), végül a tömegét. Energiát a modell nem ad meg: az étkezés energiáját a 
    \texttt{benchmark.py} a válasz feldolgozásakor számolja, tételenként a tömeg és a 100 grammra 
    vonatkozó energiasűrűség szorzatának századaként, majd ezeket összegzi. A prompt 
    ételcsoportonként jellemző energiasűrűség-tartományokat is megad. 
    Makrotápanyagokat nem becsül.
\end{enumerate}
\paragraph{}A P3\_anchor, a P3\_portion, a P3\_check és a P3\_recipe\_check a P3\_recipe-hez 
hasonlóan összes szénhidrátot kér (a rosttal együtt), a P3 adagbecslési szabályait pedig mind az 
öt változat szó szerint megtartja.
A promptok teljes szövegét a Melléklet \ref{sec:promptszovegek}.~alfejezete közli.

\paragraph{}A P3\_anchor, a P3\_portion, a P3\_check és a P3\_density azt vizsgálja, javít-e a 
becslésen, ha a modell általános tudását kifejezett lépésekre és durva viszonyítási 
tartományokra bontom. Külső tudásforrást, például tápanyag-adatbázist ezért nem vontam be, 
ellentétben a FNDDS-re támaszkodó DietAI24-gyel \cite{yan2025dietai24}. Az energia 
számításához a promptok az általános Atwater-faktorokat adják meg \cite{fao2003foodenergy}, 
a térfogatokhoz szokásos konyhai mértékegységeket (teáskanál 5~ml, evőkanál 15~ml) és a 
folyadékok nagyjából 1~g/ml sűrűségét. Az ételcsoportonkénti összetételi és energiasűrűség-tartományokat, 
a tányér-, pohár- és dobozméreteket, az ételsűrűségeket és az ellenőrzés 15\%-os tűrését a 
promptok összeállításakor általános ismeretek alapján, nagy nyelvi modell (Claude Opus 5.5) segítségével adtam meg, 
durva támpontként. A tárgyméreteket tartományként adtam meg, mert egy adott fényképen 
a tányér vagy a pohár valódi mérete nem ismert. Az ellenőrzésben azért van tűrés, mert az 
általános Atwater-faktorokkal számolt energia az egyes ételek tényleges energiatartalmától 
eltérhet \cite{fao2003foodenergy}. A 15\%-os érték választott határ.

\subsection{A bemenetek típusai}
\label{sec:bemenet_tipusok}
\paragraph{}A mérésekben négy bemenettípus különböztethető meg. A szöveges rész eredete 
típusonként eltér, ezért az eredményeik nem vonhatók össze, és az eredmények táblázataiban 
mindig megadom, melyik típusról van szó.
\begin{enumerate}
    \item \textbf{Csak kép:} egy vagy több fénykép, szöveg nélkül. Ide tartozik a P3, a P3\_macros 
    és a többi képes változat, ezekben a futásokban a modell a képen kívül semmilyen leírást nem kapott.
    \item \textbf{Kép és a referenciarekordból képzett szöveg:} a kép mellett a résztvevő által 
    naplózott ételek FNDDS-leírása, tömeg nélkül (SNAPMe, P3\_macros\_hint). 
    A szöveg a referenciaannotációból származik, ezért pontosabban nevezi meg az ételeket, mint 
    amire egy valódi felhasználó leírása képes volna.
    \item \textbf{Csak szöveg, referenciaadatokból nyelvi modellel fogalmazva:} a NutriBench 
    leírásai, amelyeket a szerzők a felmérésekben rögzített ételekből és mennyiségekből 
    GPT-4o-mini segítségével készítettek. A NutriBench v1 grammos (metric) és hétköznapi 
    mértékegységű (natural) leírásai közül a hétköznapi változatot használtam. A szöveg az ételeket és a mennyiségüket 
    is megadja, de grammok nélkül, hétköznapi mértékegységben vagy darabszámban (például „egy 
    uncia”, „egy közepes zacskó”, „egy szendvics”).
    \item \textbf{Kép, majd a modell saját leírása:} a P3\_prep\_textcalc második lépése, amely 
    csak szöveget kap, de ezt a szöveget az első lépésben ugyanaz a modell írta a képről. Emberi 
    leírás itt sincs.
\end{enumerate}


\section{A mérési eredmények összehasonlításához használható módszerek kiválasztása}

\subsection{Hibamértékek}
\label{sec:hibamertekek}

pMAE: Minden célmennyiségre (energia, szénhidrát, fehérje, zsír, tömeg) ételenként összevetésre kerül a modell becslése a referenciaértékkel. A pMAE érték az abszolút hibák összegét a referenciaértékek összegéhez viszonyítja:
\begin{equation}
  \mathrm{pMAE} = 100\cdot\frac{\sum_i |p_i - t_i|}{\sum_i t_i},
\end{equation}
ahol $p_i$ a modell becslése, $t_i$ pedig a referenciaérték az $i$-edik ételre.
A pMAE előnyösebb a feladat szempontjából, mint a MAPE mérték (ételenkénti százalékos hiba átlaga). A MAPE minden étel hibáját a saját referenciaértékéhez (a referenciaérték az adathalmazban megadott, a becslés validációjára szolgáló érték) viszonyítja, ezért a kis referenciaértékű 
ételek már kis abszolút eltérés mellett is nagy relatív hibát adnak, és ez aránytalanul befolyásolja az átlagot. A pMAE ezzel szemben az abszolút hibák összegét a 
referenciaértékek összegéhez viszonyítja, így minden étel a tápértékével arányos súllyal számít. Ez megfelel az étkezésnaplózás céljának, amely egy adott időszak összesített 
tápanyagbevitelének minél pontosabb követése.

MAE (Mean Absolute Error): az átlagos abszolút hiba, energiatartalom esetén kilokalóriában, a többi esetben grammban.

Az előjeles torzítás (bias) megengedi, hogy az ellentétes előjelű hibák kioltsák egymást, ezért ez a mennyiség nem a hiba nagyságának a vizsgálatára szolgál,
hanem azt mutatja meg, van-e a becslés során szisztematikus torzítás, konzisztens alá- vagy fölébecslés és milyen mértékben.
A torzítást a pMAE-hez hasonlóan a referenciaértékek összegéhez viszonyítva, százalékban adom meg:
\begin{equation}
  \mathrm{bias} = 100\cdot\frac{\sum_i (p_i - t_i)}{\sum_i t_i}.
\end{equation}


\section{A mérési protokoll kialakítása}
\paragraph{}Annak érdekében, hogy egy adott mérési beállítás hatása vizsgálható legyen,
kontrollálni kell, hogy miben térhetnek el az egyes mérések. Egy nem várt körülmény
változása önmagában, vagy a szándékosan változtatott, vizsgált paraméterrel együtt hatva
olyan eredményt hozhat, amely változást idéz elő a mérés pontosságában.
Ha ezeket a körülményeket nem ismerem, téves következtetésre juthatok a vizsgált
paraméter állításával kapcsolatban.

\paragraph{}Ennek elkerülése érdekében egységes mérési protokollt alakítottam ki.
Egy mérési protokoll feladata rögzíteni minél több nem vizsgált paramétert azok közül,
amelyeknek hatása van a mérésre.
A nagy nyelvi modellek nem determinisztikus működéséből adódóan azonban nincs lehetőség
minden ilyen paraméter rögzítésére: ugyanarra a bemenetre a modell eltérő választ adhat.
Ennek a véletlenszerű ingadozásnak a nagysága ugyanazon bemenetre ismételt hívásokkal becsülhető.

A konfigurációs paramétereket az egyetlen, vizsgált paraméteren kívül rögzítve, azt az egyet változtatva
összevethető méréspárokat lehet előállítani.
Páros összehasonlítások során rögzítendő paraméter a mintaelemszám ($n$), és az úgynevezett
seed érték.
Az ételek kiválasztása az adathalmazból egy véletlenszerű sorbarendezés után történik,
az első $n$ adathalmazrekord kerül bele a vizsgált halmazba.
A seed érték rögzítésével elérhető, hogy a véletlenszerű kiválasztás reprodukálhatóan
működjön, tehát az egyes mérések ugyanazokat az ételeket fogják tartalmazni, az ételek kiválasztása
így nem torzítja a vizsgált különbséget a két mérés eredménye között.
A véletlenszerű sorbarendezés biztosítja, hogy amennyiben az adathalmaz eredeti sorrendje valamilyen rendszer
szerint épül fel, annak a hatását kontrollálja a mérésekre nézve.
Például ha növekvő sorrendben lennének tömeg szerint az ételek az adathalmazban, így nem csak a kisebb tömegűek kerülnének
bele a mérésre kiválasztott részhalmazba.
A véletlen sorbarendezés alkalmazása miatt eltekintek annak vizsgálatától, hogy tartalmaznak-e valóban az adathalmazok
ilyen, sorrend-függő torzítást a validációs értékekre nézve.

\paragraph{}A képek a küldés előtt egységes előfeldolgozáson esnek át. A \texttt{benchmark.py} program a kép EXIF-metaadataiban rögzített tájolás szerint elforgatja a képet, 
hogy ha a képek nem a megjelenítési irányban vannak tárolva, abban az esetben is egységes legyen az elforgatásuk. Ezután RGB színtérbe alakítja, és a képarány megtartásával úgy 
kicsinyíti, hogy a hosszabbik oldala legfeljebb 1024 képpont legyen. Az ennél kisebb képeket nem nagyítja. Végül a képet 90-es minőségi beállítással JPEG formátumba kódolja. 
Az egységes előfeldolgozás biztosítja, hogy az eltérő forrású adathalmazok képei azonos felbontásban és formátumban jussanak el a modellekhez, és a bemeneti tokenszám, 
a költség bemeneti oldalról meghatározó része felülről korlátozott legyen.
A JPEG-tömörítés mértékének hatását páros előméréssel vizsgáltam (Gemini Flash-Lite, Nutrition5k, P3\_macros, 50 étel). A 80-as minőségi beállítással az energia pMAE-je 2,9, a 
szénhidráté 4,4 százalékponttal nagyobb volt, mint a 90-es beállítással. A különbség ekkora mintán nem kimutatható (95\%-os konfidenciaintervallum: $-3{,}2$ és $+9{,}1$, 
illetve $-3{,}7$ és $+12{,}6$ százalékpont között), de nem utal arra, hogy az erősebb tömörítés elegendő lenne. Mivel a bemeneti tokenszám a két beállítás mellett azonos volt,
a 90-es minőséget választottam.
A Pillow programkönyvtár dokumentációja a 95 feletti értékeket nem javasolja: a 100-as beállítás a tömörítés egyes lépéseit kikapcsolja, ami nagy fájlméretet eredményez 
a képminőség érdemi javulása nélkül \cite{pillow_jpeg}.

\paragraph{}Ha egy API-hívás hibával tér vissza (például hálózati hiba, túlterhelés vagy a kérések számának korlátozása miatt), a program legfeljebb háromszor újrapróbálkozik, 
vagyis hívásonként legfeljebb négy kísérletet tesz. A kísérletek között exponenciálisan növekvő ideig vár (5, 10, majd 20 másodpercig), hogy a szolgáltató átmeneti túlterhelése 
ne okozzon azonnali újabb hibát. Ha a négy kísérlet egyike sem sikeres, a hívást sikertelenként naplózza. A válaszidőt mindig csak a sikeres kísérletre méri, így a válaszidő a 
hibás kísérletek és a várakozások idejét nem tartalmazza.

\subsection{Előmérési protokoll}
\paragraph{}Az előmérési protokoll feladata, hogy minél több mérési paraméter
hatását feltárhatóvá tegye, miközben kontrollálja az ebből adódó többletköltség
nagyságát. Ennek érdekében a mintaszámot ezen mérések során a döntési mérési protokollhoz
képest kisebb értékben, adathalmazonként 50 ételben rögzítettem. Ez több összehasonlítást tesz lehetővé, viszont
kisebb pontossággal teszi megismerhetővé a vizsgált paraméter hatását. A kisebb elemszám
hatására adódhatnak olyan mérések, amelyek során a mérésre kiválasztott minták az adathalmazból
kedveznek a mért paraméter pontosságának, vagy rontják azt.
Az előmérési protokoll célja azon paraméterek hatásának feltárása, amelyek olyan
mértékben határozzák meg a mérések pontosságát, hogy ilyen elemszám mellett is képesek
kimutatható eltérést produkálni.

\subsubsection{Az előmérési protokoll által vizsgált konfigurációs paraméterek bemutatása}
Ezek azok a konfigurációs beállítások, amelyek hatását össze lehet hasonlítani
a mérési protokoll segítségével:
\begin{enumerate}
    \item Prompt utasítás:
    A becslési feladat szöveges megfogalmazása a mért nagy nyelvi modell számára.
    \item Nagy nyelvi modell:
    A vizsgált nagy nyelvi modell kiválasztására szolgál.
    \item Token keret:
    A generálható kimeneti tokenek maximális száma becslésenként, a költség limitálása
    céljából
    \item Temperature paraméter:
    Azt szabályozza, hogy a nagy nyelvi modell a
    következő token kiválasztásakor mennyire véletlenszerűen választ a lehetséges tokenek
    közül.
    \item "Reasoning" kapcsoló:
    Azt szabályozza, hogy a nagy nyelvi modell a
    végső válasz előtt mennyi belső következtetést állít elő. Ezek az úgynevezett
    érvelési tokenek a válaszban nem jelennek meg, de a modell ezekre támaszkodva
    alakítja ki a kimenetet.
    \item Kép típus:
    A DFood-Tj adathalmaz esetén felülnézeti és ferde (nagyjából 45°-os szögű), az OmniFood8k adathalmaz esetén a hat különböző szögből készült kép használatának kombinációját lehet a paraméterrel konfigurálni.
\end{enumerate}

\subsubsection{Az előmérési protokoll által mért kimeneti paraméterek bemutatása}
Ezek azok az értékek, amelyekre hatást gyakorolnak a konfigurációs beállítások,
és változásuk vizsgálhatóvá teszi, hogyan teljesített a feladatban az adott konfiguráció.
\begin{enumerate}
    \item API-hívások száma
    \item Az újrapróbálkozások után is sikertelen API-hívások száma
    \item Fel nem dolgozható válaszok száma
    \item Csak a JSON-objektum kiemelése után feldolgozható válaszok száma
    \item Átlagos válaszidő
    \item Medián (P50) válaszidő
    \item Hívásonkénti bemeneti tokenszám mediánja
    \item Költség
    \item Az energiatartalom, a tömeg, valamint a szénhidrát-, fehérje- és zsírtartalom
    becslésének hibája (pMAE, MAE, bias)
\end{enumerate}

A \texttt{benchmark.py} program minden hívásnál naplózza a bemeneti és a kimeneti tokenek számát, és ezekből a 
szolgáltató által a mérés idején közzétett listaárak alapján kiszámítja a hívás költségét dollárban. A 
méréseket összesítően leíró kimeneti fájlban a költség mellett azért szerepel a hívásonkénti bemeneti 
tokenszám mediánja, mert azt nagyrészt a mérés beállításai határozzák meg, így a konfigurációt jellemzi. 
Ez a tulajdonsága a beállított és a mért paraméterek határára helyezi: közvetlenül nem hangolható a 
mérések konfigurációja során, de azokból származó mennyiség, emellett függ a használt nagy nyelvi 
modell tokenizálási módszerétől is. Mivel közvetlenül nem állítható be, a mért kimeneti paraméterek 
közé sorolom. A kimeneti tokenszám ezzel szemben a modell válaszától függ. Ez a tokenkerettel felülről 
korlátozható, hatása pedig a költségben jelenik meg.

\subsection{Döntési mérési protokoll}
\label{sec:dontesi_protokoll}
\paragraph{}A döntési mérési protokoll célja, hogy az előmérési protokoll során feltárt eredményeket 
nagyobb, korábban nem vizsgált mintán, előre rögzített eljárással ellenőrizze, és ezzel megalapozza, 
mely beállítások kerüljenek a kialakítandó becslési folyamatba.

\paragraph{}A vizsgált modellek a Gemini 3.5 Flash-Lite, 
a Claude Sonnet 5.5 és a GPT-6.1 Sol, a \ref{tab:modellek}.~táblázatban szereplő beállításokkal, 
hívásonként egy képpel és egy ismétléssel.

\paragraph{}A minták az előmérésekkel azonos, seed szerinti véletlen sorrendből származnak, de 
nagyobb elemszámmal. A futások a sorrend elejét, vagyis az előmérésekben már felhasznált ételeket 
is tartalmazzák, ezeket azonban kizárom a hipotézisvizsgálatokból, így a tesztek egy kivétellel
korábban nem elemzett ételeken alapulnak. A kivétel a DFood-Tj adathalmazon jelentkezik: a vizsgált 250 ételéből 25 szerepelt 
az első mérési kör két próbafuttatásában, amelyek még a végső mintavételi eljárás előtt készültek, 
más prompttal, és csak az energiát és a tömeget becsülték, a szénhidrátot nem. Ezen ételek nélkül 
(225 páron) a különbségek kisebbek, de a következtetés nem változik: Sonnetre $-4{,}9$ ($-6{,}7$; $-3{,}1$), 
Solra $-7{,}0$ ($-9{,}7$; $-4{,}3$) százalékpont, mindkettő Holm-korrekció után is szignifikáns. A tesztelt minta a Nutrition5k, az 
ACETADA és a SNAPMe adathalmazon a sorrend 51--250. étele, az OmniFood8k-n a 101--300., a DFood-Tj-n a 351--600. étele. Az előmérésekben 
már felhasznált ételek eredményeit csak a konzisztencia vizsgálatára használom: ugyanezen ételek 
korábbi, azonos beállítású futásainak eredményeivel vetem össze őket.

\paragraph{}A hipotéziseket családokba soroltam (\ref{tab:hipotezisek}.~táblázat). 
Minden összehasonlítás kétoldali páros $t$-próba a pMAE különbségére, ugyanazokon az ételeken. 
Egy különbséget akkor tekintek szignifikánsnak, ha a családon belül Holm-korrekcióval javított 
$p$-érték kisebb, mint 0,05. A hatás nagyságát korrekció nélküli 95\%-os konfidenciaintervallummal 
közlöm.
Ez a becslés pontosságát mutatja, a szignifikanciáról a korrigált $p$-érték dönt.

\paragraph{}Az A család vizsgálatába azok az adathalmazok tartoznak, amelyeken az előmérésben 
talált különbség legalább az egyik drágább modellnél elérte a 200 elemes mintán kimutatható 
legkisebb különbséget (MDD). Ezt páros összehasonlításra a
\begin{equation}
  \mathrm{MDD} = \left(z_{1-\alpha'/2} + z_{1-\beta}\right)\frac{\sigma_d}{\sqrt{n}}
\end{equation}
összefüggéssel számoltam, ahol $\sigma_d$ az ételenkénti hibakülönbségek szórása az 
előmérésben, $n$ a minta elemszáma, $\alpha' = 0{,}05/m$ a Holm-korrekció legszigorúbb 
lépésének szintje egy $m$ elemű családban, az A családban $m = 8$, $1-\beta$ pedig a próba ereje. Az 
erőt 80\%-ra választottam: egy pontosan MDD nagyságú valódi 
különbséget a próba 80\% valószínűséggel mutat ki. A SNAPMe adathalmaz ennek a feltételnek nem 
felelt meg.

\paragraph{}Az előmérések nagyszámú páros összehasonlításának áttekintéséhez segédeszközként 
egy összesített pontszámot is használtam: adathalmazonként megszámoltam, hány összehasonlításban 
volt egy konfiguráció a 95\%-os konfidenciaintervallum alapján jobb vagy rosszabb a többinél, 
és ezt az adathalmazokra átlagoltam. Ez az áttekintés a döntési mérés hipotéziseinek 
kiválasztását segítette. Eredményként nem közlöm, mert az eltérő konyhájú, felvételi 
körülményű és referenciaadatú adathalmazokat egyenlő súllyal vonja össze, és sok páros 
összehasonlításon alapul, többszörös összehasonlításra vonatkozó korrekció nélkül, így a 
véletlenül kimutatottnak látszó különbségek is beleszámítanak.

\begin{table}[h]
\centering
\caption{A döntési mérés hipotéziscsaládjai. Minden kérdést kétoldali páros próbával vizsgáltam, 
a nullhipotézis szerint a két konfiguráció pMAE-je nem tér el. $m$: a családon belüli 
összehasonlítások száma.}\label{tab:hipotezisek}
\begin{tabular}{lp{5.0cm}p{3.6cm}lr}
\hline
Család & Vizsgált kérdés & Adathalmaz & Cél & $m$ \\
\hline
A  & Pontosabb-e a Flash-Lite szénhidrátban a Sonnetnél és a Solnál? (P3\_macros) & Nutrition5k, OmniFood8k, ACETADA, DFood-Tj & szénhidrát & 8 \\
B  & Javít-e a referenciarekordból képzett ételnév-lista (tömeg nélkül) a csak képes becslésen? (P3\_macros\_hint vs.\ P3\_macros, Flash-Lite) & SNAPMe & szénhidrát, zsír & 2 \\
C1 & Megismétlődik-e a P3\_recipe zsírbecslési javulása nagyobb mintán? (Flash-Lite) & OmniFood8k & zsír & 1 \\
C2 & Változtat-e a P3\_recipe a szénhidrátbecslés hibáján? (Flash-Lite) & OmniFood8k, Nutrition5k & szénhidrát & 2 \\
C3 & Javít-e a P3\_recipe a zsírbecslésen a Sol modellnél is? & OmniFood8k & zsír & 1 \\
\hline
\end{tabular}
\end{table}

\paragraph{}A kiértékelés szabályait a futtatás előtt rögzítettem:
kiugró értékek szerinti kizárást nem alkalmazok, a fel nem dolgozható 
válaszokat nem kérem le újra, és a pontossági eredményeket csak azután 
vizsgálom, hogy egy hipotéziscsalád minden mérése lefutott. 
A hipotézisekhez nem tartozó mennyiségeket teszt nélkül, 
95\%-os konfidenciaintervallummal írom le.

\subsection{A vizsgálati kérdések és a mérések kapcsolata}
\label{sec:kerdes_meres}
\paragraph{}\Az{\ref{tab:kerdes_meres}}.~táblázat a bevezetésben megfogalmazott kérdésekhez rendeli 
a megválaszolásukhoz használt méréseket. Döntési szempontot csak azoknál a kérdéseknél adok meg, amelyeket a döntési 
mérés előre rögzített hipotézisként vizsgált (\ref{tab:hipotezisek}.~táblázat). A többi kérdésre a válasz leíró: 
az eredményeket 95\%-os konfidenciaintervallummal közlöm, szignifikanciaállítás nélkül.
A 3. és a 4. kérdés két, külön mérésekkel vizsgált részből áll. Ezeket a táblázat a és b betűvel különbözteti meg.

\begin{table}[htbp]
\centering
\footnotesize
\setlength{\tabcolsep}{3pt}
\caption{A vizsgálati kérdések és a hozzájuk tartozó mérések. A „Döntési szempont” oszlopban 
a „Holm” azt jelenti, hogy egy összehasonlítás akkor szignifikáns, ha a hipotéziscsaládon belül 
Holm-korrigált $p$-értéke $0{,}05$ alatt van.}
\label{tab:kerdes_meres}
\begin{tabular}{>{\raggedright\arraybackslash}p{1.0cm}
                >{\raggedright\arraybackslash}p{3.0cm}
                >{\raggedright\arraybackslash}p{2.7cm}
                >{\raggedright\arraybackslash}p{2.5cm}
                >{\raggedright\arraybackslash}p{2.2cm}
                >{\raggedright\arraybackslash}p{2.0cm}}
\hline
Kérdés & Bemenet és konfiguráció & Adathalmaz & Elsődleges mérőszám & Összehasonlítás & Döntési szempont \\
\hline
1. & kép, P3\_macros; előmérés: 8 modell, adathalmazonként 50 étel; döntési mérés: Flash-Lite, Sonnet, Sol 
   & előmérés: Nutrition5k, OmniFood8k, ACETADA, DFood-Tj, SNAPMe; döntési mérés: ugyanezek a SNAPMe nélkül 
   & szénhidrát pMAE; kiegészítő: energia, tömeg, fehérje, zsír 
   & páros, adathalmazonként 
   & A család ($m = 8$), Holm; a többi leíró \\
2. & kép, P3\_macros, Flash-Lite 
   & amelyekben a tömeg is ismert 
   & az energiahiba tömeg- és energiasűrűség-része 
   & \hspace{0pt}adathalmazonként 
   & leíró \\
3a. & 1 kép vagy több nézőpont, P3 vagy P3\_macros, Flash-Lite 
   & DFood-Tj (2 nézet), OmniFood8k (6 kamera) 
   & energia pMAE 
   & páros: 1 kép és több kép 
   & leíró \\
3b. & kép, promptváltozatok, főként Flash-Lite 
   & előmérés: változónként; döntési mérés: OmniFood8k, Nutrition5k 
   & energia pMAE; döntési mérés: zsír, szénhidrát 
   & páros, a P3\_macros-hoz képest 
   & C1--C3 család, Holm; a többi leíró \\
4a. & csak szöveg, NutriBench prompt, 7 modell 
   & NutriBench 
   & energia, szénhidrát pMAE 
   & modellek között 
   & leíró \\
4b. & kép és a referenciarekordból képzett ételnév-lista, Flash-Lite 
   & SNAPMe 
   & szénhidrát, zsír pMAE 
   & páros, a csak képes P3\_macros-hoz képest 
   & B család ($m = 2$), Holm; felső becslés \\
5. & a döntési mérés futásai; konzisztencia: két független hívás 
   & a döntési mérés adathalmazai 
   & válaszidő, költség étkezésenként, sikeres hívások aránya; ételenkénti szórás 
   & modellek között, adathalmazonként 
   & leíró \\
6. & kép, P3\_macros tétellistái 
   & Nutrition5k, DFood-Tj, ACETADA, SNAPMe 
   & lefedettség és találati arány a referencia-tételnevekhez képest 
   & megnevezett és címkézett tételek 
   & leíró \\
\hline
\end{tabular}
\end{table}

\section{Az utólagos kalibráció elvetése}
Az első mérési körben a Gemini Flash-Lite becslései szisztematikus torzítást mutattak: 
a nagyobb adagokat a kisebbekhez képest alulbecsülte, és az energiaértékek átlagos előjeles hibája 
adathalmazonként erősen eltért. Ez a torzítás
az ACETADA adathalmazon elhanyagolható volt (\mbox{$-1{,}6\%$}), a Nutrition5k-n
mérsékelt (\mbox{$-11{,}1\%$}), az OmniFood8k-n viszont nagy (\mbox{$-45{,}0\%$}).
Egy adott adathalmazon illesztett utólagos korrekció, például a becsült és a valós érték
közötti lineáris leképezés, ugyanazon az adathalmazon képes csökkenteni ezt a rendszeres
eltérést: az első mérési körben a kihagyásos (leave-one-out) módon illesztett lineáris
korrekció az OmniFood8k-n az energia pMAE-jét 50,4\%-ról 32,1\%-ra csökkentette.
Viszont amellett döntöttem, hogy ezt az elkészítendő becslést végző pipeline-ba nem fogom
beépíteni.
\paragraph{}A döntés oka, hogy a korrekció csak abban a helyzetben érvényes, amelyből
illesztették. A torzítás nagysága és előjele függ az adathalmaztól, az ételek jellegétől, a használt modelltől és az adag méretétől,
ezért egyetlen rögzített korrekciós függvény az egyik helyzetben javítana, a másikban
viszont ronthatna. A mérések ezt alátámasztják: ugyanez a korrekció a Nutrition5k-n
(31,7\%-ról 34,1\%-ra) és az ACETADA-n (22,6\%-ról 23,6\%-ra) rontott az eredményen.
Egy valós alkalmazásban a bemenetek
eloszlása ismeretlen és változó, a kalibrációhoz szükséges címkézett minták pedig nem
állnak rendelkezésre, így a korrekció nem lenne megbízhatóan beállítható.


\clearpage % Ez azért kell, hogy nehogy képek átcsússzanak a következő fejezethez
